% ===================================================================================
% BAB 1: PENDAHULUAN
% Berisi pengenalan umum mengenai metode pemilihan OSIS dan aplikasi e-SPPO.
% File ini akan di-include oleh main.tex
% ===================================================================================

\part{Pendahuluan}
\label{part:pendahuluan}

Dokumen ini merupakan panduan resmi untuk penggunaan dan administrasi \textbf{Sistem Penyelenggaraan Pemilihan OSIS Elektronik (e-SPPO)} versi 2.0.0. Panduan ini dirancang untuk menjadi referensi utama bagi seluruh pihak yang terlibat dalam pemilihan, termasuk panitia penyelenggara, pengawas pemilihan, dan para pemilih. Tujuannya adalah untuk memastikan bahwa seluruh proses pemilihan dapat berjalan dengan lancar, transparan, dan sesuai dengan prosedur yang ditetapkan.

% -----------------------------------------------------------------------------------
% Section: Metode-Metode Pelaksanaan Pemilihan OSIS
% -----------------------------------------------------------------------------------
\section{Metode-Metode Pelaksanaan Pemilihan OSIS}
\label{sec:metode_pemilihan}

Seiring perkembangan zaman, pelaksanaan pemilihan Ketua dan Wakil Ketua OSIS telah berevolusi. Pemilihan metode yang tepat sering kali bergantung pada skala pemilihan, sumber daya, dan infrastruktur yang dimiliki oleh sekolah. Secara umum, terdapat dua pendekatan utama: konvensional dan modern.

\subsection{Metode Konvensional (Manual)}
Metode ini adalah pendekatan tradisional yang mengandalkan proses fisik sepenuhnya. Pemilih menggunakan surat suara kertas untuk mencoblos atau menandai pasangan calon pilihan mereka, yang kemudian dimasukkan ke dalam kotak suara. Setelah periode pemungutan suara berakhir, panitia akan membuka kotak suara dan menghitung setiap suara secara manual di hadapan saksi.

\begin{itemize}
	\item \textbf{Kelebihan:} Prosesnya sangat sederhana dan mudah dipahami oleh semua orang. Tidak memerlukan investasi teknologi yang signifikan, sehingga dapat diterapkan di sekolah dengan keterbatasan fasilitas.
	\item \textbf{Kekurangan:} Proses penghitungan suara memakan waktu yang sangat lama, terutama jika jumlah pemilih besar. Metode ini juga sangat rentan terhadap kesalahan manusia (\textit{human error}), baik dalam proses pencoblosan (suara tidak sah) maupun saat rekapitulasi. Selain itu, penggunaan kertas dalam jumlah besar berdampak buruk bagi lingkungan.
\end{itemize}

\subsection{Metode Modern (Elektronik)}
Metode modern memanfaatkan teknologi untuk mengatasi kelemahan dari metode konvensional, terutama dalam hal kecepatan, akurasi, dan efisiensi.

\subsubsection{Metode Semi-Elektronik}
Metode ini merupakan jembatan antara sistem manual dan digital penuh. Pemilih tetap menggunakan surat suara fisik, namun didesain dalam format khusus seperti \textbf{Lembar Jawaban Komputer (LJK)}. Setelah pemungutan suara, seluruh surat suara akan dipindai (\textit{scanned}) menggunakan mesin pemindai khusus. Perangkat lunak kemudian akan membaca tanda pada LJK dan menghitung hasilnya secara otomatis. Ini secara signifikan mengurangi waktu rekapitulasi dan meminimalisir kesalahan hitung.

\subsubsection{Metode Elektronik Penuh (\textit{E-Voting})}
Metode ini sepenuhnya meninggalkan penggunaan kertas. Seluruh proses, mulai dari pemungutan hingga penghitungan suara, dilakukan secara digital. Pemilih mengakses sebuah aplikasi atau sistem, yang disebut \textbf{Surat Suara Elektronik (SSE)}, melalui perangkat seperti komputer, tablet, atau ponsel pintar. Setelah melakukan autentikasi, pemilih dapat langsung memilih kandidat pada layar. Suara yang masuk akan direkam dan dihitung oleh server secara \textit{real-time}. \textbf{Sistem e-SPPO termasuk dalam kategori ini}, menyediakan solusi \textit{e-voting} yang terintegrasi.

% -----------------------------------------------------------------------------------
% Section: Apa Itu e-SPPO?
% -----------------------------------------------------------------------------------
\section{Apa Itu e-SPPO?}
\label{sec:apa_itu_esppo}

\textbf{Sistem Penyelenggaraan Pemilihan OSIS Elektronik (e-SPPO)} adalah sebuah \textit{software suite} (kumpulan program aplikasi) berbasis web yang komprehensif, dirancang untuk mengelola dan melaksanakan seluruh tahapan pemilihan Ketua dan Wakil Ketua OSIS secara digital.

e-SPPO bukan sekadar aplikasi pemungutan suara, melainkan sebuah ekosistem terpadu yang mencakup manajemen data pemilih, pendaftaran kandidat, proses pemungutan suara yang aman, hingga rekapitulasi dan pelaporan hasil secara otomatis. Sistem ini lahir dari kebutuhan akan sebuah platform pemilihan yang tidak hanya efisien, tetapi juga menjunjung tinggi nilai-nilai transparansi, akuntabilitas, dan keamanan data dalam lingkungan pendidikan.

% -----------------------------------------------------------------------------------
% Section: Fungsi e-SPPO
% -----------------------------------------------------------------------------------
\section{Fungsi e-SPPO}
\label{sec:fungsi_esppo}
Sebagai sebuah sistem terpadu, e-SPPO dirancang untuk menjalankan empat fungsi krusial dalam penyelenggaraan pemilihan OSIS:
\begin{itemize}
	\item \textbf{Pengelolaan Data Terpusat:} e-SPPO menyediakan panel administrasi bagi panitia untuk mengelola seluruh data pemilihan di satu tempat. Ini mencakup impor dan validasi Daftar Pemilih Tetap (DPT), pendaftaran pasangan calon beserta visi, misi, dan foto, serta manajemen akun untuk panitia dan pengawas.
	\item \textbf{Pelaksanaan Pemungutan Suara Elektronik:} Fungsi inti dari e-SPPO adalah menyediakan antarmuka Surat Suara Elektronik (SSE) yang intuitif dan aman bagi pemilih. Sistem memastikan bahwa setiap pemilih yang sah hanya dapat memberikan satu suara, dan kerahasiaan pilihan mereka tetap terjaga.
	\item \textbf{Rekapitulasi Suara Otomatis dan Real-time:} Setiap suara yang masuk melalui SSE akan secara otomatis direkam dan diakumulasikan oleh sistem. Panitia dan pengawas dapat memantau perolehan suara secara langsung (\textit{live count}) melalui dasbor administrasi, menghilangkan kebutuhan akan penghitungan manual yang melelahkan dan berisiko.
	\item \textbf{Pemantauan dan Pelaporan Komprehensif:} Sistem ini memberikan alat bagi panitia untuk memantau tingkat partisipasi pemilih. Setelah pemilihan selesai, e-SPPO dapat secara otomatis menghasilkan berbagai dokumen pelaporan, seperti berita acara hasil pemilihan dan rekapitulasi suara total, yang siap untuk dicetak dan ditandatangani.
\end{itemize}

% -----------------------------------------------------------------------------------
% Section: Sejarah Metode Pemilihan OSIS di SMABAT dan e-SPPO
% -----------------------------------------------------------------------------------
\section{Sejarah Metode Pemilihan OSIS di SMABAT dan e-SPPO}
\label{sec:sejarah_esppo}
e-SPPO adalah puncak dari serangkaian inovasi dan perbaikan sistem pemilihan yang telah diterapkan di SMAN 1 Bati-Bati. Setiap tahap evolusi didasari oleh evaluasi terhadap kekurangan sistem sebelumnya dan keinginan untuk menciptakan proses yang lebih baik.
\begin{enumerate}
	\item \textbf{Metode Manual (1992 -- pertengahan 2010-an):} Sejak pendirian sekolah, pemilihan dilakukan secara konvensional. Awalnya, proses ini dilakukan dengan cara panitia membawa kotak suara berkeliling dari kelas ke kelas. Pada pertengahan dekade 2010-an, sistem ini beralih ke model Tempat Pemungutan Suara (TPS), di mana pemilih datang ke lokasi tertentu untuk memberikan suara. Di akhir periode, seluruh surat suara dihitung secara terpusat. Meskipun prosesnya berjalan, tantangan utama dari metode manual ini adalah waktu dan tenaga yang besar untuk rekapitulasi serta potensi suara tidak sah.
	\item \textbf{Elektronik dengan Google Form (2020-2021):} Pandemi COVID-19 memaksa diterapkannya pembelajaran jarak jauh, yang secara langsung berdampak pada metode pemilihan. Google Form digunakan sebagai solusi darurat untuk memungkinkan pemilihan tetap berlangsung secara elektronik. Namun, metode ini memiliki kelemahan signifikan dalam hal keamanan, validasi pemilih (risiko \textit{double voting}), dan kerahasiaan suara.
	\item \textbf{Aplikasi E-Pilketos (2022-2023):} Seiring implementasi Kurikulum Merdeka, semangat Proyek Penguatan Profil Pelajar Pancasila (P5) dengan tema "Suara Demokrasi" mendorong adopsi teknologi \textit{e-voting}. Aplikasi \textit{open-source} E-Pilketos pertama kali digunakan pada tahun 2022 sebagai tugas proyek P5 oleh kelas X-D dan X-C. Aplikasi ini menjadi demonstrasi teknologi \textit{e-voting} dan pembeda dengan kelas lain yang masih menggunakan metode manual. Atas keberhasilan proyek tersebut, sistem ini kemudian diadopsi secara resmi oleh OSIS dan digunakan untuk pemilihan tahun 2022 dan 2023. Namun, E-Pilketos memiliki beberapa kelemahan fundamental: dikembangkan oleh pihak luar, menggunakan \textit{framework} CodeIgniter versi lama yang tidak kompatibel dengan PHP 8.x, menyimpan data pemilih secara \textit{plain-text} (terbuka), dan menggunakan \textit{packages} yang usang.
	\item \textbf{e-SPPO (2024 - Sekarang):} Untuk mengatasi kelemahan E-Pilketos, e-SPPO dikembangkan secara mandiri (\textit{in-house}) sebagai penggantinya. Versi 1.0.0, yang digunakan untuk pemilihan tahun 2024, membawa perbaikan signifikan: dibangun menggunakan kode PHP native yang kompatibel dengan versi 8.x, memanfaatkan Composer untuk manajemen \textit{packages}, dan mengenkripsi data pemilih dalam suara yang tersimpan menggunakan \texttt{password\_hash()}. Meskipun versi ini masih memiliki banyak kekurangan, hasilnya terbukti jauh lebih memuaskan. Keberhasilan tersebut menjadi fondasi pengembangan versi 2.0.0, yang hadir dengan perbaikan arsitektur, keamanan, dan fitur yang lebih matang.
\end{enumerate}

% -----------------------------------------------------------------------------------
% Section: Kenapa e-SPPO? (Keunggulan)
% -----------------------------------------------------------------------------------
\section{Kenapa e-SPPO? (Keunggulan)}
\label{sec:keunggulan_esppo}
e-SPPO menawarkan serangkaian keunggulan yang menjadikannya solusi ideal untuk pemilihan OSIS di era digital:
\begin{itemize}
	\item \textbf{Keamanan Data Lanjutan:} Keamanan adalah prioritas utama. Pada versi 1.0.0, data pemilih dalam suara dienkripsi menggunakan \texttt{password\_hash()}, yang memberikan peningkatan keamanan secara \textbf{marginal} karena keterbatasan fungsi \textit{hash} Bcrypt untuk tujuan enkripsi. Pada versi 2.0.0, sistem ini beralih ke enkripsi simetris \textbf{AES-256-CBC} yang jauh lebih kuat untuk melindungi data pemilih. Sementara itu, untuk autentikasi akun, sistem tetap mempertahankan mekanisme \textit{hashing} kata sandi yang aman menggunakan \texttt{password\_hash()} dan \texttt{password\_verify()}.
	\item \textbf{Kemudahan Akses dan Penggunaan:} Dibangun dengan \textit{framework} antarmuka modern seperti \textbf{Bootstrap 5} dan \textbf{AdminLTE 3.2}, e-SPPO menjamin pengalaman pengguna yang optimal. Tampilannya responsif, artinya dapat diakses dengan nyaman dari berbagai perangkat, mulai dari PC desktop hingga tablet dan ponsel pintar.
	\item \textbf{Transparansi dan Akuntabilitas:} e-SPPO menjamin bahwa setiap suara terekam secara akurat dan dihitung secara tepat waktu. Selama pemilihan berlangsung, hasil hanya dapat diakses oleh pihak berwenang. Namun, setelah pemilihan selesai, hasil dapat dipublikasikan melalui portal publik untuk transparansi penuh. Sistem ini juga mampu menghasilkan beragam laporan dalam format PDF untuk keperluan dokumentasi dan diseminasi informasi.
	\item \textbf{Fleksibilitas Implementasi:} e-SPPO dirancang untuk sangat adaptif. Sistem ini mendukung pemisahan antara server web dan server basis data, serta dapat diimplementasikan di berbagai lingkungan jaringan: mulai dari intranet yang terisolasi penuh (\textit{air-gapped network}), jaringan lokal dengan akses internet terbatas, hingga di-hosting di \textit{cloud}. Metode pelaksanaannya pun fleksibel, dapat disesuaikan dengan kebutuhan, baik berbasis TPS dengan perangkat khusus (PASSE) di setiap bilik suara maupun diakses secara jarak jauh oleh pemilih.
\end{itemize}

% -----------------------------------------------------------------------------------
% Section: Arsitektur Dasar dan Komponen e-SPPO
% -----------------------------------------------------------------------------------
\section{Arsitektur Dasar dan Komponen e-SPPO}
\label{sec:arsitektur_esppo}
Sebagai aplikasi berbasis web, e-SPPO berjalan di atas arsitektur \textit{client-server}. Pengguna (klien) mengakses aplikasi melalui peramban web, sementara seluruh logika pemrosesan dan penyimpanan data terjadi di sisi server. Infrastruktur server ini dapat diatur dalam dua konfigurasi utama: \textbf{tunggal (\textit{stand-alone})} di mana server web dan basis data berjalan di satu mesin, atau \textbf{terpisah (\textit{split server})} di mana keduanya berjalan di mesin yang berbeda untuk skalabilitas yang lebih baik.

Sistem ini mendukung akses melalui protokol \textbf{HTTP} maupun \textbf{HTTPS}, tergantung pada konfigurasi server web. Penggunaan HTTPS sangat diutamakan, terutama jika sistem dapat diakses dari jaringan internet, untuk menjamin enkripsi data selama transit. Untuk jaringan lokal yang aman dan terisolasi, HTTP dapat menjadi pilihan.

\subsection{Teknologi yang Digunakan}
Fondasi teknologi e-SPPO dipilih untuk memaksimalkan fleksibilitas, keamanan, dan kemudahan pemeliharaan:
\begin{itemize}
	\item \textbf{Back-End:} Menggunakan \textbf{PHP murni} tanpa \textit{framework} tambahan. Pilihan ini memberikan kontrol penuh atas kode, performa yang ringan, dan kemudahan untuk di-deploy di hampir semua layanan hosting web standar.
	\item \textbf{Front-End:} Kombinasi \textit{framework} antarmuka yang berbeda digunakan untuk setiap komponen. \textbf{Bootstrap 5} menjadi dasar untuk \textit{Surat Suara Elektronik (SSE)} dan \textit{Portal Publik}, memberikan tampilan yang bersih dan modern. Sementara itu, \textbf{AdminLTE 3.2} digunakan untuk \textit{Pusminlihdu}, menyediakan templat dasbor yang kaya fitur dan ideal untuk administrasi data.
	\item \textbf{Basis Data:} Sistem berkomunikasi menggunakan protokol \textbf{MySQL}, sehingga kompatibel dengan server basis data \textbf{MySQL} maupun \textbf{MariaDB}, dua sistem manajemen basis data relasional yang paling populer dan andal.
\end{itemize}

\subsection{Komponen Utama e-SPPO Versi 2.0.0}
Secara fungsional, ekosistem e-SPPO terbagi menjadi dua komponen aplikasi utama dan satu komponen pendukung yang saling terintegrasi:

\subsubsection{Surat Suara Elektronik (SSE)}
Ini adalah "wajah" dari e-SPPO yang dilihat oleh \textbf{pemilih}. SSE adalah aplikasi web yang menjadi gerbang bagi pemilih untuk menggunakan hak suaranya. Alur penggunaannya dirancang agar lugas dan sederhana: pemilih membuka aplikasi, melakukan \textit{login} dengan akun yang telah diberikan, memilih salah satu kandidat, mengonfirmasi pilihan tersebut untuk dikirim dan diproses oleh sistem, kemudian melakukan \textit{logout} untuk mengakhiri sesi. Antarmuka SSE sengaja dibuat minimalis untuk memastikan semua pemilih dapat menggunakannya dengan mudah tanpa kebingungan.

\subsubsection{Pusat Administrasi Pemilihan Terpadu (Pusminlihdu)}
Pusminlihdu adalah "ruang kontrol" atau pusat komando dari seluruh kegiatan pemilihan. Ini adalah panel administrasi komprehensif yang hanya dapat diakses oleh \textbf{Panitia Seleksi dan Pemilihan (Panselih)} serta \textbf{Pengawas Pemilihan} dengan hak akses berjenjang. Dari sini, panitia dapat mengelola berbagai jenis data dan konfigurasi, antara lain:
\begin{itemize}
	\item \textbf{Data Administratif:} Mengelola akun untuk pengelola (administrator) dan pengawas, serta memasukkan data pejabat kesiswaan yang akan menandatangani dokumen resmi.
	\item \textbf{Detail Kegiatan Pemilihan:} Mengatur parameter krusial seperti nama kegiatan, jadwal (tanggal mulai dan selesai), jumlah TPS, jumlah Perangkat Akses Surat Suara Elektronik (PASSE), tipe kandidat (tunggal atau berpasangan), mode tampilan SSE, hingga status pemilihan (aktif, selesai, dll).
	\item \textbf{Data Konstituensi dan DPT:} Mengelola data pemilih, mulai dari pengelompokan (misalnya per kelas) hingga Daftar Pemilih Tetap (DPT) yang akan digunakan dalam pemilihan.
	\item \textbf{Data Calon:} Memasukkan dan mengelola informasi mengenai pasangan calon, termasuk nama, foto, visi, dan misi.
	\item \textbf{Fungsi Lainnya:} Memantau jalannya pemilihan secara \textit{real-time}, melihat hasil perolehan suara, dan yang terpenting, membuat berbagai laporan dan berkas pemilihan dalam format PDF.
\end{itemize}

\subsubsection{Portal Web Publik e-SPPO}
Ini adalah komponen pendukung yang berfungsi sebagai etalase informasi pemilihan kepada publik atau seluruh warga sekolah. Portal ini dirancang untuk menampilkan hasil akhir rekapitulasi suara secara transparan setelah seluruh proses pemilihan selesai dan diverifikasi oleh panitia. Selain itu, portal ini juga menjadi tempat di mana laporan dan berkas-berkas lain yang dihasilkan oleh Pusminlihdu dapat diunduh oleh umum, sehingga memperkuat akuntabilitas proses pemilihan.

% --- Akhir dari Bab 1: Pendahuluan ---
